\begin{figure}[htbp]
\centering
\begin{tikzpicture}[
    font=\footnotesize,
    >=Stealth,
    baseline/.style={draw=statsblue!55, line width=0.65pt},
    signal/.style={draw=statsviolet!75!black, thick, -{Stealth[length=2.1mm]}},
    panel/.style={rounded corners=7pt, draw=black!18, fill=black!1},
    info/.style={
        rectangle, rounded corners=4pt, align=center,
        minimum width=35mm, minimum height=11mm, font=\scriptsize
    },
    note/.style={align=center, font=\scriptsize}
]

\path[use as bounding box] (-7.15,-4.05) rectangle (7.15,3.75);

% -----------------------------------------------------------------
% Campo interferométrico
% -----------------------------------------------------------------
\fill[statsmint!3, rounded corners=8pt] (-7.0,-3.35) rectangle (2.35,3.25);
\draw[statsmint!35!black, rounded corners=8pt, line width=0.7pt]
    (-7.0,-3.35) rectangle (2.35,3.25);
\node[font=\small\bfseries, text=statsnavy] at (-2.32,3.52)
    {ARREGLO COMPLETO: $n=20$ ANTENAS};

% Posiciones inspiradas en brazos de una espiral logarítmica.
\foreach \i/\x/\y in {
  1/-2.40/0.00, 2/-1.75/0.28, 3/-1.38/0.92,
  4/-1.62/1.72, 5/-2.45/2.18, 6/-3.55/2.05,
  7/-4.42/1.38, 8/-4.72/0.35, 9/-4.38/-0.80,
  10/-3.52/-1.72, 11/-2.30/-2.15, 12/-0.92/-1.88,
  13/0.10/-0.98, 14/0.55/0.28, 15/0.25/1.62,
  16/-0.62/2.55, 17/-2.12/2.78, 18/-3.92/2.58,
  19/-5.55/1.48, 20/-5.95/-0.20}
  \coordinate (ant\i) at (\x,\y);

% Zona con restricción: cuatro antenas afectadas por sombra orográfica.
\fill[statspink!7, rounded corners=12pt]
    (-6.35,0.72) rectangle (-3.05,3.00);
\draw[statspink!65, dashed, thick, rounded corners=12pt]
    (-6.35,0.72) rectangle (-3.05,3.00);
\node[fill=white, inner sep=1.5pt, text=statspink!70!black,
      font=\scriptsize\bfseries] at (-4.70,2.78)
    {ZONA DE SOMBRA};
\node[
  note, text=statspink!70!black, fill=white,
  rounded corners=2pt, inner sep=1.5pt
] at (-4.70,0.96)
    {4 antenas; 2 activas\\(máximo permitido)};

% Las quince líneas de base del sub-array activo de seis antenas.
\foreach \a/\b in {
  2/6,2/8,2/11,2/14,2/19,
  6/8,6/11,6/14,6/19,
  8/11,8/14,8/19,
  11/14,11/19,
  14/19}
  \draw[baseline] (ant\a) -- (ant\b);

% Antenas no seleccionadas.
\foreach \i in {1,3,4,5,7,9,10,12,13,15,16,17,18,20}{
  \begin{scope}[shift={(ant\i)}]
    \fill[white] (0,0) circle (0.18);
    \draw[black!45, line width=0.65pt] (0,0) circle (0.18);
    \draw[black!50, line width=0.7pt] (-0.12,0.08)
      arc[start angle=150,end angle=30,radius=0.14];
    \draw[black!50, line width=0.7pt] (0,-0.02) -- (0,-0.23);
    \draw[black!50, line width=0.7pt] (-0.12,-0.23) -- (0.12,-0.23);
  \end{scope}
}

% Seis antenas seleccionadas para el sub-array.
\foreach \i in {2,6,8,11,14,19}{
  \begin{scope}[shift={(ant\i)}]
    \fill[statsgold!92] (0,0) circle (0.23);
    \draw[statsnavy, line width=0.9pt] (0,0) circle (0.23);
    \draw[statsnavy, line width=0.85pt] (-0.14,0.10)
      arc[start angle=150,end angle=30,radius=0.17];
    \draw[statsnavy, line width=0.85pt] (0,-0.02) -- (0,-0.28);
    \draw[statsnavy, line width=0.85pt] (-0.14,-0.28) -- (0.14,-0.28);
    \node[font=\tiny\bfseries, text=statsnavy, above=2.5mm] {$A_{\i}$};
  \end{scope}
}

\node[
  rectangle, rounded corners=4pt, draw=statsblue!70!black,
  fill=statsblue!7, inner sep=3pt, align=center,
  font=\scriptsize\bfseries
] at (-2.30,-2.88)
  {Sub-array activo: $k=6$\\$\binom{6}{2}=15$ líneas de base};

% -----------------------------------------------------------------
% Correlador y significado de las líneas de base
% -----------------------------------------------------------------
\draw[signal] (2.02,0.05) -- (3.05,0.05);
\node[font=\tiny, text=statsviolet!75!black, above] at (2.54,0.05)
    {señales};

\fill[statsnavy!4, rounded corners=8pt] (3.10,-3.35) rectangle (7.0,3.25);
\draw[statsnavy!40, rounded corners=8pt, line width=0.7pt]
    (3.10,-3.35) rectangle (7.0,3.25);
\node[font=\small\bfseries, text=statsnavy] at (5.05,3.52)
    {CORRELADOR DIGITAL};

\node[
  rectangle, rounded corners=5pt, draw=statsviolet!70!black,
  fill=statsviolet!8, minimum width=31mm, minimum height=14mm,
  align=center, font=\bfseries
] (corr) at (5.05,1.73)
  {Combina cada par\\de antenas activas};

\node[info, draw=statsblue!70!black, fill=statsblue!7]
  (pairs) at (5.05,0.28)
  {$15$ pares no ordenados\\$A_iA_j$, con $i<j$};

\node[info, draw=statsmint!65!black, fill=statsmint!8]
  (uv) at (5.05,-1.12)
  {Muestreo del plano $uv$\\por cada línea de base};

\node[info, draw=statspink!65!black, fill=statspink!7]
  (image) at (5.05,-2.52)
  {Síntesis de apertura\\y resolución angular};

\draw[signal] (corr) -- (pairs);
\draw[signal] (pairs) -- (uv);
\draw[signal] (uv) -- (image);

% Leyenda compacta.
\fill[statsgold!92] (3.52,-3.72) circle (0.10);
\node[anchor=west, font=\scriptsize] at (3.70,-3.72) {antena activa};
\draw[baseline] (5.05,-3.72) -- (5.55,-3.72);
\node[anchor=west, font=\scriptsize] at (5.66,-3.72) {línea de base};

\end{tikzpicture}
\caption{Representación de la configuración de un sub-array en radiointerferometría. De las 20 antenas disponibles se seleccionan 6, resaltadas en amarillo. Cada par de antenas activas forma una línea de base, por lo que el sub-array genera $\binom{6}{2}=15$ líneas de base que el correlador utiliza para sintetizar una apertura. La región rosada contiene las cuatro antenas sometidas a la restricción de sombra orográfica; en el ejemplo se seleccionan exactamente dos, el máximo permitido.}
\end{figure}